\documentclass[11pt,a4paper]{article}
\usepackage[utf8]{inputenc}
\usepackage[T1]{fontenc}
\usepackage[spanish,es-minimal,es-noshorthands]{babel}
\AtBeginDocument{\renewcommand{\sin}{\operatorname{sen}}\renewcommand{\arcsin}{\operatorname{arcsen}}}
\usepackage{lmodern}
\usepackage[sfdefault,lining]{FiraSans}
\usepackage[a4paper,top=2.3cm,bottom=2.2cm,left=1.8cm,right=1.8cm,headheight=1.3cm,headsep=0.4cm,footskip=1cm]{geometry}
\usepackage{amsmath,amssymb}
\usepackage{icomma}
\usepackage{graphicx}
\usepackage[table]{xcolor}
\usepackage{tikz}
\usepackage{fontawesome5}
\usepackage{ccicons}
\usepackage{fancyhdr}
\usepackage{lastpage}
\usepackage{tasks}
\usepackage{needspace}
\usepackage[most]{tcolorbox}
\usepackage{array}
\usepackage[hidelinks]{hyperref}
\AtBeginDocument{\hypersetup{pdftitle={Ficha de ejercicios de la unidad 10: Estadística. 1º Bach.},pdfsubject={Matemáticas I 1º Bach.},pdfkeywords={estadística, frecuencias, media, mediana, cuartiles, varianza, desviación típica, coeficiente de variación, variable bidimensional, tabla de doble entrada, nube de puntos, covarianza, correlación lineal}}}

\definecolor{navy}{HTML}{1B2748}
\definecolor{accent}{HTML}{E8603A}
\definecolor{teal}{HTML}{1F9E90}
\definecolor{soft}{HTML}{EEF1F8}
\definecolor{muted}{HTML}{5D6685}

\setlength{\parindent}{0pt}
\setlength{\parskip}{3pt}
\settasks{label=\alph*), label-format=\bfseries\color{navy}, label-width=1.3em, item-indent=1.8em, column-sep=0.6em, after-item-skip=5pt, before-skip=3pt, after-skip=4pt}

% ---------- capçalera i peu ----------
\newcommand{\llicencia}{\href{https://creativecommons.org/licenses/by-nc-sa/4.0/deed.es}{\ccbyncsa}\ \ Licencia Creative Commons Atribución-NoComercial-CompartirIgual 4.0 Internacional (CC BY-NC-SA 4.0)}
\pagestyle{fancy}
\fancyhf{}
\fancyhead[L]{\raisebox{-0.35\height}{\includegraphics[height=0.95cm]{logo-gva.png}}}
\fancyhead[C]{\small\color{muted}Ficha de ejercicios · Unidad 10: Estadística · 1\textsuperscript{o} Bach.}
\fancyhead[R]{\raisebox{-0.35\height}{\includegraphics[height=0.95cm]{logo-vermellar.png}}}
\fancyfoot[L]{\scriptsize\color{muted}\llicencia}
\fancyfoot[R]{\small\color{muted}\thepage}
\renewcommand{\headrulewidth}{0.4pt}
\renewcommand{\headrule}{\color{muted!40}\hrule height \headrulewidth}
\fancypagestyle{primera}{\fancyhf{}\renewcommand{\headrulewidth}{0pt}
  \fancyfoot[L]{\scriptsize\color{muted}\llicencia}
  \fancyfoot[R]{\small\color{muted}\thepage}}

% ---------- nivells i exercicis ----------
\newcommand{\estrelles}[1]{{\color{accent}\ifcase#1\or\faStar\faStar[regular]\faStar[regular]\or\faStar\faStar\faStar[regular]\or\faStar\faStar\faStar\fi}}
\newcounter{exer}
\newcommand{\ex}[1]{\par\needspace{4\baselineskip}\medskip\refstepcounter{exer}\noindent{\color{navy}\bfseries\theexer.}\ \ }

\newcommand{\seccio}[2]{\par\needspace{16\baselineskip}\bigskip
  \begin{tcolorbox}[colback=navy,colframe=navy,coltext=white,arc=2mm,boxrule=0pt,left=3mm,right=3mm,top=1.2mm,bottom=1.2mm]
  {\large\bfseries #1\quad #2}
  \end{tcolorbox}\vspace{-2pt}}

\newcommand{\exf}[1]{\par\medskip\refstepcounter{exer}\noindent{\color{navy}\bfseries\theexer.}\ \ }
\providecommand{\gt}{>}\providecommand{\lt}{<}
% quadre de recordatori (només a la fitxa d'exercicis)
\newcommand{\recorda}[1]{\par\smallskip\begin{tcolorbox}[colback=accent!7,colframe=accent!55,arc=2mm,boxrule=0.5pt,left=2.5mm,right=2.5mm,top=1mm,bottom=1mm]\small{\color{accent}\faLightbulb[regular]}\ \textbf{Recuerda:}\ #1\end{tcolorbox}}
\newcommand{\res}[1]{#1}\newcommand{\rest}[1]{#1}

\def\fmneed{11}% que el títol d'una secció no quede sol al final del full
% =====================================================================
% Capçalera FentMates de les fitxes: un color per a cada curs.
% Els colors són els de la pàgina inicial (LEVELS a assets/js/materials.js);
% el logo ve de _eines/logo/pdf/ (python3 _eines/logo/genera.py).
%
% Ús, al final del preàmbul de la unitat:
%   % =====================================================================
% Capçalera FentMates de les fitxes: un color per a cada curs.
% Els colors són els de la pàgina inicial (LEVELS a assets/js/materials.js);
% el logo ve de _eines/logo/pdf/ (python3 _eines/logo/genera.py).
%
% Ús, al final del preàmbul de la unitat:
%   % =====================================================================
% Capçalera FentMates de les fitxes: un color per a cada curs.
% Els colors són els de la pàgina inicial (LEVELS a assets/js/materials.js);
% el logo ve de _eines/logo/pdf/ (python3 _eines/logo/genera.py).
%
% Ús, al final del preàmbul de la unitat:
%   \input{../comu/fentmates.tex}
%   \fmcurs{1eso}
% i, al començament del document:
%   \capcalera{Matemàtiques}{1r ESO}{Fitxa d'exercicis · Unitat 2}{Divisibilitat}
% Opcional abans de \input: \def\fmneed{14} (espai mínim abans de cada secció).
% =====================================================================
\makeatletter
\def\fm@logodir{../../logo/pdf}
\providecommand{\fmneed}{10}
\def\fm@cursnom{}

% clau del curs, color i color del text damunt del color
\newcommand{\fm@curs}[3]{\@namedef{fm@color@#1}{#2}\@namedef{fm@text@#1}{#3}}
\fm@curs{1eso}{E6B84A}{navy}
\fm@curs{2eso}{3B82F6}{white}
\fm@curs{3eso}{FF7A50}{white}
\fm@curs{4eso-a}{A78BFA}{white}
\fm@curs{4eso-b}{6D28D9}{white}
\fm@curs{1batx-mi}{9B1C1C}{white}
\fm@curs{1batx-ccss}{EF6B6B}{white}
\fm@curs{2batx-mii}{1E6B47}{white}
\fm@curs{2batx-ccss}{2BB3A3}{white}

\newcommand{\fmcurs}[1]{%
  \@ifundefined{fm@color@#1}{\PackageError{fentmates}{Curs desconegut: #1}{}}{}%
  \edef\fm@tmp{\noexpand\definecolor{curs}{HTML}{\@nameuse{fm@color@#1}}%
               \noexpand\colorlet{curstext}{\@nameuse{fm@text@#1}}}\fm@tmp
  \colorlet{cursfosc}{curs!72!black}%
  \def\fm@logo{\fm@logodir/vertical-#1.pdf}%
  \def\fm@marca{\fm@logodir/marca-#1.pdf}}

% ---------- capçalera de la primera pàgina ----------
% #1 matèria (petit), #2 curs, #3 tipus de document i unitat, #4 títol
\newcommand{\capcalera}[4]{%
  \gdef\fm@cursnom{#2}%
  \par\noindent
  \begin{tcolorbox}[colback=curs,colframe=curs,boxrule=0pt,arc=2.2mm,width=3.5cm,height=2.15cm,
      valign=center,left=3mm,right=2mm,top=1mm,bottom=1mm,nobeforeafter,box align=center]
    \color{curstext}{\footnotesize #1}\par\vspace{1pt}{\fontsize{19}{21}\selectfont\bfseries #2\par}%
  \end{tcolorbox}\hspace{5mm}%
  \begin{minipage}[c]{\dimexpr\linewidth-3.5cm-5mm-2.9cm\relax}
    {\color{cursfosc}\small\bfseries\MakeUppercase{#3}}\par\vspace{2pt}
    {\color{navy}\fontsize{20}{23}\selectfont\bfseries #4\par}
  \end{minipage}\hfill
  \begin{minipage}[c]{2.6cm}\raggedleft\includegraphics[width=2.6cm]{\fm@logo}\end{minipage}\par
  \vspace{2pt}}

% ---------- barres de secció del color del curs ----------
\renewcommand{\seccio}[2]{\par\needspace{\fmneed\baselineskip}\bigskip
  \begin{tcolorbox}[colback=curs,colframe=curs,coltext=curstext,arc=2mm,boxrule=0pt,left=3mm,right=3mm,top=1.2mm,bottom=1.2mm]
  {\large\bfseries #1\quad #2}
  \end{tcolorbox}\vspace{-2pt}}

% ---------- estrelles de dificultat del color del curs (a Batxillerat no se'n posen) ----------
\@ifundefined{estrelles}{}{%
  \renewcommand{\estrelles}[1]{{\color{curs}\ifcase#1\or\faStar\faStar[regular]\faStar[regular]\or\faStar\faStar\faStar[regular]\or\faStar\faStar\faStar\fi}}}

% ---------- capçalera i peu de les pàgines ----------
% primera pàgina: logos institucionals al peu
\fancypagestyle{primera}{\fancyhf{}\renewcommand{\headrulewidth}{0pt}
  \fancyfoot[L]{\raisebox{-0.32\height}{\includegraphics[height=0.78cm]{logo-gva.png}}\hspace{5mm}%
    \raisebox{-0.32\height}{\includegraphics[height=0.78cm]{logo-vermellar.png}}}
  \fancyfoot[R]{\scriptsize\color{muted}\href{https://creativecommons.org/licenses/by-nc-sa/4.0/deed.ca}{\ccbyncsa}\ \ CC BY-NC-SA 4.0 · fentmates.org\qquad\small\thepage}}
% la resta: el símbol a l'esquerra i el curs a la dreta, amb el títol al centre com abans
\fancyhead[L]{\raisebox{-0.3\height}{\includegraphics[height=0.62cm]{\fm@marca}}}
\fancyhead[R]{\small\bfseries\color{cursfosc}\fm@cursnom}
\renewcommand{\headrule}{\color{curs!45}\hrule height \headrulewidth}
\makeatother

%   \fmcurs{1eso}
% i, al començament del document:
%   \capcalera{Matemàtiques}{1r ESO}{Fitxa d'exercicis · Unitat 2}{Divisibilitat}
% Opcional abans de \input: \def\fmneed{14} (espai mínim abans de cada secció).
% =====================================================================
\makeatletter
\def\fm@logodir{../../logo/pdf}
\providecommand{\fmneed}{10}
\def\fm@cursnom{}

% clau del curs, color i color del text damunt del color
\newcommand{\fm@curs}[3]{\@namedef{fm@color@#1}{#2}\@namedef{fm@text@#1}{#3}}
\fm@curs{1eso}{E6B84A}{navy}
\fm@curs{2eso}{3B82F6}{white}
\fm@curs{3eso}{FF7A50}{white}
\fm@curs{4eso-a}{A78BFA}{white}
\fm@curs{4eso-b}{6D28D9}{white}
\fm@curs{1batx-mi}{9B1C1C}{white}
\fm@curs{1batx-ccss}{EF6B6B}{white}
\fm@curs{2batx-mii}{1E6B47}{white}
\fm@curs{2batx-ccss}{2BB3A3}{white}

\newcommand{\fmcurs}[1]{%
  \@ifundefined{fm@color@#1}{\PackageError{fentmates}{Curs desconegut: #1}{}}{}%
  \edef\fm@tmp{\noexpand\definecolor{curs}{HTML}{\@nameuse{fm@color@#1}}%
               \noexpand\colorlet{curstext}{\@nameuse{fm@text@#1}}}\fm@tmp
  \colorlet{cursfosc}{curs!72!black}%
  \def\fm@logo{\fm@logodir/vertical-#1.pdf}%
  \def\fm@marca{\fm@logodir/marca-#1.pdf}}

% ---------- capçalera de la primera pàgina ----------
% #1 matèria (petit), #2 curs, #3 tipus de document i unitat, #4 títol
\newcommand{\capcalera}[4]{%
  \gdef\fm@cursnom{#2}%
  \par\noindent
  \begin{tcolorbox}[colback=curs,colframe=curs,boxrule=0pt,arc=2.2mm,width=3.5cm,height=2.15cm,
      valign=center,left=3mm,right=2mm,top=1mm,bottom=1mm,nobeforeafter,box align=center]
    \color{curstext}{\footnotesize #1}\par\vspace{1pt}{\fontsize{19}{21}\selectfont\bfseries #2\par}%
  \end{tcolorbox}\hspace{5mm}%
  \begin{minipage}[c]{\dimexpr\linewidth-3.5cm-5mm-2.9cm\relax}
    {\color{cursfosc}\small\bfseries\MakeUppercase{#3}}\par\vspace{2pt}
    {\color{navy}\fontsize{20}{23}\selectfont\bfseries #4\par}
  \end{minipage}\hfill
  \begin{minipage}[c]{2.6cm}\raggedleft\includegraphics[width=2.6cm]{\fm@logo}\end{minipage}\par
  \vspace{2pt}}

% ---------- barres de secció del color del curs ----------
\renewcommand{\seccio}[2]{\par\needspace{\fmneed\baselineskip}\bigskip
  \begin{tcolorbox}[colback=curs,colframe=curs,coltext=curstext,arc=2mm,boxrule=0pt,left=3mm,right=3mm,top=1.2mm,bottom=1.2mm]
  {\large\bfseries #1\quad #2}
  \end{tcolorbox}\vspace{-2pt}}

% ---------- estrelles de dificultat del color del curs (a Batxillerat no se'n posen) ----------
\@ifundefined{estrelles}{}{%
  \renewcommand{\estrelles}[1]{{\color{curs}\ifcase#1\or\faStar\faStar[regular]\faStar[regular]\or\faStar\faStar\faStar[regular]\or\faStar\faStar\faStar\fi}}}

% ---------- capçalera i peu de les pàgines ----------
% primera pàgina: logos institucionals al peu
\fancypagestyle{primera}{\fancyhf{}\renewcommand{\headrulewidth}{0pt}
  \fancyfoot[L]{\raisebox{-0.32\height}{\includegraphics[height=0.78cm]{logo-gva.png}}\hspace{5mm}%
    \raisebox{-0.32\height}{\includegraphics[height=0.78cm]{logo-vermellar.png}}}
  \fancyfoot[R]{\scriptsize\color{muted}\href{https://creativecommons.org/licenses/by-nc-sa/4.0/deed.ca}{\ccbyncsa}\ \ CC BY-NC-SA 4.0 · fentmates.org\qquad\small\thepage}}
% la resta: el símbol a l'esquerra i el curs a la dreta, amb el títol al centre com abans
\fancyhead[L]{\raisebox{-0.3\height}{\includegraphics[height=0.62cm]{\fm@marca}}}
\fancyhead[R]{\small\bfseries\color{cursfosc}\fm@cursnom}
\renewcommand{\headrule}{\color{curs!45}\hrule height \headrulewidth}
\makeatother

%   \fmcurs{1eso}
% i, al començament del document:
%   \capcalera{Matemàtiques}{1r ESO}{Fitxa d'exercicis · Unitat 2}{Divisibilitat}
% Opcional abans de \input: \def\fmneed{14} (espai mínim abans de cada secció).
% =====================================================================
\makeatletter
\def\fm@logodir{../../logo/pdf}
\providecommand{\fmneed}{10}
\def\fm@cursnom{}

% clau del curs, color i color del text damunt del color
\newcommand{\fm@curs}[3]{\@namedef{fm@color@#1}{#2}\@namedef{fm@text@#1}{#3}}
\fm@curs{1eso}{E6B84A}{navy}
\fm@curs{2eso}{3B82F6}{white}
\fm@curs{3eso}{FF7A50}{white}
\fm@curs{4eso-a}{A78BFA}{white}
\fm@curs{4eso-b}{6D28D9}{white}
\fm@curs{1batx-mi}{9B1C1C}{white}
\fm@curs{1batx-ccss}{EF6B6B}{white}
\fm@curs{2batx-mii}{1E6B47}{white}
\fm@curs{2batx-ccss}{2BB3A3}{white}

\newcommand{\fmcurs}[1]{%
  \@ifundefined{fm@color@#1}{\PackageError{fentmates}{Curs desconegut: #1}{}}{}%
  \edef\fm@tmp{\noexpand\definecolor{curs}{HTML}{\@nameuse{fm@color@#1}}%
               \noexpand\colorlet{curstext}{\@nameuse{fm@text@#1}}}\fm@tmp
  \colorlet{cursfosc}{curs!72!black}%
  \def\fm@logo{\fm@logodir/vertical-#1.pdf}%
  \def\fm@marca{\fm@logodir/marca-#1.pdf}}

% ---------- capçalera de la primera pàgina ----------
% #1 matèria (petit), #2 curs, #3 tipus de document i unitat, #4 títol
\newcommand{\capcalera}[4]{%
  \gdef\fm@cursnom{#2}%
  \par\noindent
  \begin{tcolorbox}[colback=curs,colframe=curs,boxrule=0pt,arc=2.2mm,width=3.5cm,height=2.15cm,
      valign=center,left=3mm,right=2mm,top=1mm,bottom=1mm,nobeforeafter,box align=center]
    \color{curstext}{\footnotesize #1}\par\vspace{1pt}{\fontsize{19}{21}\selectfont\bfseries #2\par}%
  \end{tcolorbox}\hspace{5mm}%
  \begin{minipage}[c]{\dimexpr\linewidth-3.5cm-5mm-2.9cm\relax}
    {\color{cursfosc}\small\bfseries\MakeUppercase{#3}}\par\vspace{2pt}
    {\color{navy}\fontsize{20}{23}\selectfont\bfseries #4\par}
  \end{minipage}\hfill
  \begin{minipage}[c]{2.6cm}\raggedleft\includegraphics[width=2.6cm]{\fm@logo}\end{minipage}\par
  \vspace{2pt}}

% ---------- barres de secció del color del curs ----------
\renewcommand{\seccio}[2]{\par\needspace{\fmneed\baselineskip}\bigskip
  \begin{tcolorbox}[colback=curs,colframe=curs,coltext=curstext,arc=2mm,boxrule=0pt,left=3mm,right=3mm,top=1.2mm,bottom=1.2mm]
  {\large\bfseries #1\quad #2}
  \end{tcolorbox}\vspace{-2pt}}

% ---------- estrelles de dificultat del color del curs (a Batxillerat no se'n posen) ----------
\@ifundefined{estrelles}{}{%
  \renewcommand{\estrelles}[1]{{\color{curs}\ifcase#1\or\faStar\faStar[regular]\faStar[regular]\or\faStar\faStar\faStar[regular]\or\faStar\faStar\faStar\fi}}}

% ---------- capçalera i peu de les pàgines ----------
% primera pàgina: logos institucionals al peu
\fancypagestyle{primera}{\fancyhf{}\renewcommand{\headrulewidth}{0pt}
  \fancyfoot[L]{\raisebox{-0.32\height}{\includegraphics[height=0.78cm]{logo-gva.png}}\hspace{5mm}%
    \raisebox{-0.32\height}{\includegraphics[height=0.78cm]{logo-vermellar.png}}}
  \fancyfoot[R]{\scriptsize\color{muted}\href{https://creativecommons.org/licenses/by-nc-sa/4.0/deed.ca}{\ccbyncsa}\ \ CC BY-NC-SA 4.0 · fentmates.org\qquad\small\thepage}}
% la resta: el símbol a l'esquerra i el curs a la dreta, amb el títol al centre com abans
\fancyhead[L]{\raisebox{-0.3\height}{\includegraphics[height=0.62cm]{\fm@marca}}}
\fancyhead[R]{\small\bfseries\color{cursfosc}\fm@cursnom}
\renewcommand{\headrule}{\color{curs!45}\hrule height \headrulewidth}
\makeatother

\fmcurs{1batx-mi}

\begin{document}
\thispagestyle{primera}

\capcalera{Matemáticas I}{1\textsuperscript{o} Bach.}{Ficha de ejercicios · Unidad 10}{\raggedright Estadística}
\vspace{10pt}
Nombre y apellidos: \rule{6.9cm}{0.4pt}\hfill Grupo: \rule{1.6cm}{0.4pt}\hfill Fecha: \rule{2.2cm}{0.4pt}
\vspace{8pt}

\begin{tcolorbox}[colback=soft,colframe=soft,arc=2mm,left=3mm,right=3mm,top=1.5mm,bottom=1.5mm]
\small Resuelve los ejercicios en el cuaderno mostrando todos los pasos: para las variables bidimensionales, haz siempre la tabla con las columnas $x_i$, $y_i$, $x_i^2$, $y_i^2$ y $x_iy_i$. Redondea los resultados a dos o tres cifras decimales. Puedes comprobarlos con las calculadoras de los apuntes interactivos.
\end{tcolorbox}

% =====================================================
\seccio{1}{Estadística unidimensional: variables y tablas}

\exf{1} Clasifica las variables estadísticas siguientes en cualitativas, cuantitativas discretas o cuantitativas continuas:\begin{tasks}(2)
  \task El grupo sanguíneo.
  \task El número de libros leídos en un año.
  \task La temperatura a mediodía.
  \task El tiempo semanal dedicado al deporte.
  \task La marca del móvil.
  \task El número de goles en un partido.
\end{tasks}

\ex{1} Las notas de un examen de 25 alumnos han sido: \ $5$, $7$, $6$, $4$, $8$, $6$, $5$, $7$, $6$, $9$, $3$, $6$, $7$, $5$, $6$, $8$, $4$, $7$, $6$, $5$, $6$, $8$, $7$, $6$, $5$. \ Haz la tabla de frecuencias absolutas, relativas y acumuladas. ¿Qué porcentaje de alumnos ha aprobado (nota 5 o más)?\recorda{$h_i=\frac{f_i}{N}$; \ las frecuencias acumuladas $F_i$ y $H_i$ suman las de los valores menores o iguales que $x_i$.}


% =====================================================
\seccio{2}{Parámetros de centralización y de posición}

\exf{1} Calcula la media, la mediana, la moda y los cuartiles:\begin{tasks}(1)
  \task $2,\ 4,\ 4,\ 5,\ 6,\ 7,\ 7,\ 7,\ 9$
  \task $12,\ 15,\ 15,\ 18,\ 20,\ 21,\ 25,\ 30$
\end{tasks}\recorda{$\bar x=\frac{\sum x_if_i}{N}$. La mediana deja la mitad de los datos por debajo; $Q_1$, el 25\,\% y $Q_3$, el 75\,\%. Si $\frac N2$ cae justo entre dos datos diferentes, la mediana es la media de ambos.}

\ex{1} Con la tabla del ejercicio 2, calcula la nota media, la mediana, la moda y los cuartiles.

\ex{2} La media de 4 notas de una alumna es 6,5. ¿Qué nota tiene que sacar en el quinto examen para que la media sea 7?


% =====================================================
\seccio{3}{Parámetros de dispersión}

\exf{1} Dos grupos han sacado estas notas. Calcula la media, el rango, la varianza y la desviación típica de cada grupo. ¿Qué observas?\begin{tasks}(1)
  \task Grupo A: $6,\ 7,\ 7,\ 8,\ 8,\ 8,\ 9,\ 9,\ 10$
  \task Grupo B: $2,\ 4,\ 6,\ 8,\ 8,\ 8,\ 10,\ 12,\ 14$
\end{tasks}\recorda{$\sigma^2=\frac{\sum x_i^2f_i}{N}-\bar x^2$, \ $\sigma=\sqrt{\sigma^2}$, \ $CV=\frac{\sigma}{\bar x}$.}

\ex{2} Los minutos que 40 alumnos dedican a leer cada día se recogen en esta tabla. Calcula la media, la desviación típica y el coeficiente de variación.\par\smallskip{\centering\small$\begin{array}{|c|ccccc|}\hline\text{Minutos}&[0,10)&[10,20)&[20,30)&[30,40)&[40,50)\\\hline f_i&6&10&15&7&2\\\hline\end{array}$\par}

\ex{2} Dos atletas han hecho estos tiempos (s) en cinco carreras de 100 m. Calcula la media y la desviación típica de cada uno. ¿Cuál es más rápido de media? ¿Cuál es más regular?\begin{tasks}(1)
  \task A: $12{,}1,\ 12{,}3,\ 12,\ 12{,}4,\ 12{,}2$
  \task B: $11{,}8,\ 12{,}6,\ 12,\ 12{,}9,\ 11{,}7$
\end{tasks}

\ex{2} Las alturas de un grupo tienen $\bar x=165$ cm y $\sigma=8$ cm, y los pesos de otro grupo, $\bar x=32$ kg y $\sigma=4$ kg. ¿Cuál de los dos es relativamente más disperso? ¿Por qué no se pueden comparar directamente las desviaciones típicas?


% =====================================================
\seccio{4}{Variables bidimensionales. Tablas de doble entrada}

\exf{1} Se ha preguntado a 25 familias el número de hijos ($X$) y el número de coches ($Y$) que tienen:\par\smallskip{\centering {\small$\begin{array}{|c|c|c|c|}\hline X\backslash Y & 0 & 1 & 2\\\hline 1 & 4 & 2 & 0\\ 2 & 3 & 6 & 3\\ 3 & 0 & 3 & 4\\\hline \end{array}$}\par}\smallskip\begin{tasks}(1)
  \task Añade las distribuciones marginales a la tabla.
  \task ¿Cuántas familias tienen 2 hijos y 1 coche? ¿Cuántas tienen más de un coche?
  \task Calcula las medias marginales $\bar x$ y $\bar y$.
\end{tasks}\recorda{En una tabla de doble entrada, las sumas de las filas dan la distribución marginal de $X$ y las de las columnas, la de $Y$. La suma de todas las casillas es $N$.}

\ex{2} Las notas de Matemáticas ($X$) y de Física ($Y$) de 12 alumnos, en una escala de 1 a 3 (1 = insuficiente, 2 = aprobado, 3 = notable o más), han sido: \ $(2,1)$, $(3,2)$, $(2,2)$, $(1,1)$, $(3,3)$, $(2,1)$, $(1,1)$, $(2,2)$, $(3,2)$, $(2,2)$, $(1,2)$, $(3,3)$. \ Construye la tabla de doble entrada y calcula las medias marginales.


% =====================================================
\seccio{5}{Nube de puntos}

\exf{1} Un coche gasta gasolina y anotamos el combustible que queda en el depósito ($Y$, en litros) al final de cada hora de viaje ($X$): \par\smallskip{\centering {\small$\begin{array}{|c|cccccccc|}\hline x_i & 1 & 2 & 3 & 4 & 5 & 6 & 7 & 8\\\hline y_i & 62 & 60 & 55 & 53 & 50 & 46 & 44 & 40\\\hline\end{array}$}\par}\smallskip Representa la nube de puntos y el centro de gravedad. ¿Qué tipo de correlación se observa?\recorda{Cada par $(x_i,y_i)$ es un punto. El centro de gravedad es $(\bar x,\bar y)$.}

\ex{1} Asocia cada nube con uno de estos coeficientes de correlación: \ $r=0$, \ $r=0{,}9$, \ $r=-0{,}5$, \ $r=-0{,}95$.\begin{tasks}(4)
  \task \par\begin{tikzpicture}[scale=0.22,baseline=(current bounding box.center)]\draw[muted!30,very thin] (0,0) grid (10.0,10.0);\draw[-latex,thick] (0,0)--(10.5,0) node[right,font=\scriptsize]{$X$};\draw[-latex,thick] (0,0)--(0,10.5) node[above,font=\scriptsize]{$Y$};\node[below,font=\tiny] at (0.0,0) {$0$};\node[below,font=\tiny] at (2.0,0) {$2$};\node[below,font=\tiny] at (4.0,0) {$4$};\node[below,font=\tiny] at (6.0,0) {$6$};\node[below,font=\tiny] at (8.0,0) {$8$};\node[below,font=\tiny] at (10.0,0) {$10$};\node[left,font=\tiny] at (0,0.0) {$0$};\node[left,font=\tiny] at (0,2.0) {$2$};\node[left,font=\tiny] at (0,4.0) {$4$};\node[left,font=\tiny] at (0,6.0) {$6$};\node[left,font=\tiny] at (0,8.0) {$8$};\node[left,font=\tiny] at (0,10.0) {$10$};\fill[accent] (6.968,7.394) circle (0.16);\fill[accent] (7.235,7.358) circle (0.16);\fill[accent] (4.95,4.95) circle (0.16);\fill[accent] (3.577,2.806) circle (0.16);\fill[accent] (3.036,3.567) circle (0.16);\fill[accent] (4.892,4.859) circle (0.16);\fill[accent] (3.152,3.886) circle (0.16);\fill[accent] (2.467,2.897) circle (0.16);\fill[accent] (5.169,5.907) circle (0.16);\fill[accent] (5.061,4.9) circle (0.16);\fill[accent] (5.743,5.687) circle (0.16);\fill[accent] (3.33,3.996) circle (0.16);\fill[accent] (4.848,3.889) circle (0.16);\fill[accent] (4.733,4.34) circle (0.16);\fill[accent] (2.353,2.225) circle (0.16);\fill[accent] (5.729,7.1) circle (0.16);\fill[accent] (5.37,5.132) circle (0.16);\fill[accent] (8.787,8.651) circle (0.16);\fill[accent] (5.175,5.536) circle (0.16);\fill[accent] (4.601,4.324) circle (0.16);\fill[accent] (6.877,5.252) circle (0.16);\fill[accent] (5.169,5.807) circle (0.16);\fill[accent] (6.342,5.711) circle (0.16);\fill[accent] (4.236,4.812) circle (0.16);\fill[accent] (5.201,4.015) circle (0.16);\end{tikzpicture}
  \task \par\begin{tikzpicture}[scale=0.22,baseline=(current bounding box.center)]\draw[muted!30,very thin] (0,0) grid (10.0,10.0);\draw[-latex,thick] (0,0)--(10.5,0) node[right,font=\scriptsize]{$X$};\draw[-latex,thick] (0,0)--(0,10.5) node[above,font=\scriptsize]{$Y$};\node[below,font=\tiny] at (0.0,0) {$0$};\node[below,font=\tiny] at (2.0,0) {$2$};\node[below,font=\tiny] at (4.0,0) {$4$};\node[below,font=\tiny] at (6.0,0) {$6$};\node[below,font=\tiny] at (8.0,0) {$8$};\node[below,font=\tiny] at (10.0,0) {$10$};\node[left,font=\tiny] at (0,0.0) {$0$};\node[left,font=\tiny] at (0,2.0) {$2$};\node[left,font=\tiny] at (0,4.0) {$4$};\node[left,font=\tiny] at (0,6.0) {$6$};\node[left,font=\tiny] at (0,8.0) {$8$};\node[left,font=\tiny] at (0,10.0) {$10$};\fill[accent] (8.913,1.832) circle (0.16);\fill[accent] (4.499,5.153) circle (0.16);\fill[accent] (6.055,4.023) circle (0.16);\fill[accent] (5.69,4.471) circle (0.16);\fill[accent] (6.702,3.546) circle (0.16);\fill[accent] (3.412,6.621) circle (0.16);\fill[accent] (4.864,5.466) circle (0.16);\fill[accent] (4.369,4.369) circle (0.16);\fill[accent] (3.894,6.15) circle (0.16);\fill[accent] (4.233,5.793) circle (0.16);\fill[accent] (4.72,5.205) circle (0.16);\fill[accent] (5.052,5.634) circle (0.16);\fill[accent] (4.141,5.569) circle (0.16);\fill[accent] (6.095,3.981) circle (0.16);\fill[accent] (4.669,4.638) circle (0.16);\fill[accent] (0.771,9.384) circle (0.16);\fill[accent] (7.225,2.267) circle (0.16);\fill[accent] (4.898,4.589) circle (0.16);\fill[accent] (4.381,6.62) circle (0.16);\fill[accent] (5.869,4.509) circle (0.16);\fill[accent] (5.812,4.288) circle (0.16);\fill[accent] (5.552,4.616) circle (0.16);\fill[accent] (4.217,5.309) circle (0.16);\fill[accent] (5.756,3.938) circle (0.16);\fill[accent] (3.213,7.027) circle (0.16);\end{tikzpicture}
  \task \par\begin{tikzpicture}[scale=0.22,baseline=(current bounding box.center)]\draw[muted!30,very thin] (0,0) grid (10.0,10.0);\draw[-latex,thick] (0,0)--(10.5,0) node[right,font=\scriptsize]{$X$};\draw[-latex,thick] (0,0)--(0,10.5) node[above,font=\scriptsize]{$Y$};\node[below,font=\tiny] at (0.0,0) {$0$};\node[below,font=\tiny] at (2.0,0) {$2$};\node[below,font=\tiny] at (4.0,0) {$4$};\node[below,font=\tiny] at (6.0,0) {$6$};\node[below,font=\tiny] at (8.0,0) {$8$};\node[below,font=\tiny] at (10.0,0) {$10$};\node[left,font=\tiny] at (0,0.0) {$0$};\node[left,font=\tiny] at (0,2.0) {$2$};\node[left,font=\tiny] at (0,4.0) {$4$};\node[left,font=\tiny] at (0,6.0) {$6$};\node[left,font=\tiny] at (0,8.0) {$8$};\node[left,font=\tiny] at (0,10.0) {$10$};\fill[accent] (5.371,5.94) circle (0.16);\fill[accent] (7.362,6.727) circle (0.16);\fill[accent] (3.602,4.307) circle (0.16);\fill[accent] (6.918,5.012) circle (0.16);\fill[accent] (4.761,2.585) circle (0.16);\fill[accent] (4.757,4.83) circle (0.16);\fill[accent] (8.482,2.363) circle (0.16);\fill[accent] (5.479,3.479) circle (0.16);\fill[accent] (5.134,7.219) circle (0.16);\fill[accent] (6.465,2.33) circle (0.16);\fill[accent] (7.15,6.795) circle (0.16);\fill[accent] (5.154,6.07) circle (0.16);\fill[accent] (6.221,5.133) circle (0.16);\fill[accent] (3.529,6.245) circle (0.16);\fill[accent] (4.575,6.359) circle (0.16);\fill[accent] (4.452,7.127) circle (0.16);\fill[accent] (2.911,5.213) circle (0.16);\fill[accent] (2.608,4.671) circle (0.16);\fill[accent] (2.404,4.65) circle (0.16);\fill[accent] (4.796,4.114) circle (0.16);\fill[accent] (4.91,5.513) circle (0.16);\fill[accent] (4.655,4.41) circle (0.16);\fill[accent] (5.327,7.173) circle (0.16);\fill[accent] (2.906,2.778) circle (0.16);\fill[accent] (5.071,3.958) circle (0.16);\end{tikzpicture}
  \task \par\begin{tikzpicture}[scale=0.22,baseline=(current bounding box.center)]\draw[muted!30,very thin] (0,0) grid (10.0,10.0);\draw[-latex,thick] (0,0)--(10.5,0) node[right,font=\scriptsize]{$X$};\draw[-latex,thick] (0,0)--(0,10.5) node[above,font=\scriptsize]{$Y$};\node[below,font=\tiny] at (0.0,0) {$0$};\node[below,font=\tiny] at (2.0,0) {$2$};\node[below,font=\tiny] at (4.0,0) {$4$};\node[below,font=\tiny] at (6.0,0) {$6$};\node[below,font=\tiny] at (8.0,0) {$8$};\node[below,font=\tiny] at (10.0,0) {$10$};\node[left,font=\tiny] at (0,0.0) {$0$};\node[left,font=\tiny] at (0,2.0) {$2$};\node[left,font=\tiny] at (0,4.0) {$4$};\node[left,font=\tiny] at (0,6.0) {$6$};\node[left,font=\tiny] at (0,8.0) {$8$};\node[left,font=\tiny] at (0,10.0) {$10$};\fill[accent] (4.938,4.719) circle (0.16);\fill[accent] (5.633,3.133) circle (0.16);\fill[accent] (4.115,3.9) circle (0.16);\fill[accent] (5.449,4.061) circle (0.16);\fill[accent] (6.39,4.807) circle (0.16);\fill[accent] (5.546,4.358) circle (0.16);\fill[accent] (7.433,4.713) circle (0.16);\fill[accent] (3.42,8.505) circle (0.16);\fill[accent] (4.982,3.89) circle (0.16);\fill[accent] (3.715,4.178) circle (0.16);\fill[accent] (3.586,5.108) circle (0.16);\fill[accent] (4.57,7.04) circle (0.16);\fill[accent] (5.234,3.982) circle (0.16);\fill[accent] (5.558,7.351) circle (0.16);\fill[accent] (5.704,2.901) circle (0.16);\fill[accent] (8.534,2.675) circle (0.16);\fill[accent] (6.286,4.782) circle (0.16);\fill[accent] (2.256,4.433) circle (0.16);\fill[accent] (5.206,4.13) circle (0.16);\fill[accent] (3.849,6.173) circle (0.16);\fill[accent] (4.025,6.571) circle (0.16);\fill[accent] (7.013,4.881) circle (0.16);\fill[accent] (4.505,4.849) circle (0.16);\fill[accent] (1.671,8.127) circle (0.16);\fill[accent] (5.38,5.732) circle (0.16);\end{tikzpicture}
\end{tasks}

\ex{1} Indica, para cada par de variables, si esperas una correlación positiva, negativa o ninguna correlación:\begin{tasks}(1)
  \task La altura y el peso de 30 adolescentes.
  \task El precio de una entrada de cine y el número de espectadores.
  \task La nota de inglés y el número de zapato.
  \task Los kilómetros recorridos y la gasolina gastada.
  \task La temperatura exterior y el consumo de calefacción.
  \task El mes de nacimiento y el número de hermanos.
\end{tasks}


% =====================================================
\seccio{6}{Varianzas y covarianza}

\exf{2} Para la variable bidimensional siguiente, haz la tabla de cálculos y halla las medias, las varianzas, las desviaciones típicas y la covarianza:\par\smallskip{\centering {\small$\begin{array}{|c|ccccc|}\hline x_i & 2 & 4 & 5 & 6 & 8\\\hline y_i & 3 & 5 & 4 & 7 & 9\\\hline\end{array}$}\par}\recorda{$\sigma_x^2=\frac{\sum x_i^2}{N}-\bar x^2$, \ $\sigma_y^2=\frac{\sum y_i^2}{N}-\bar y^2$, \ $\sigma_{xy}=\frac{\sum x_iy_i}{N}-\bar x\,\bar y$.}

\ex{2} Calcula la covarianza de los datos siguientes e interpreta su signo:\par\smallskip{\centering {\small$\begin{array}{|c|ccccc|}\hline x_i & 10 & 20 & 30 & 40 & 50\\\hline y_i & 9 & 8 & 6 & 5 & 2\\\hline\end{array}$}\par}

\ex{2} Calcula la covarianza de las notas de Matemáticas y de Física de la sección 4 (la tabla de doble entrada de 12 alumnos).

\ex{2} Calcula la covarianza de la variable bidimensional del ejercicio sobre hijos y coches (sección 4).


% =====================================================
\seccio{7}{Correlación lineal}

\exf{2} Calcula el coeficiente de correlación de los datos de los ejercicios de la sección 6 e interprétalo.\begin{tasks}(1)
  \task Los datos del primer ejercicio de la sección 6.
  \task Los datos del segundo ejercicio de la sección 6.
  \task Hijos y coches (sección 4).
\end{tasks}\recorda{$r=\dfrac{\sigma_{xy}}{\sigma_x\,\sigma_y}$, \ $-1\leq r\leq 1$. El signo da el sentido de la relación y $|r|$, la fuerza: cerca de 1, fuerte; cerca de 0, no hay correlación lineal.}

\ex{2} La altura (cm) y el peso (kg) de 7 personas han sido:\par\smallskip{\centering {\small$\begin{array}{|c|ccccccc|}\hline x_i\ (\text{cm}) & 150 & 155 & 160 & 165 & 170 & 175 & 180\\\hline y_i\ (\text{kg}) & 48 & 52 & 55 & 61 & 63 & 70 & 72\\\hline\end{array}$}\par}\smallskip Representa la nube de puntos, calcula el coeficiente de correlación e interprétalo.

\ex{2} Las horas semanales de videojuegos ($X$) y la nota de un examen ($Y$) de 6 alumnos han sido:\par\smallskip{\centering {\small$\begin{array}{|c|cccccc|}\hline x_i & 0 & 1 & 2 & 3 & 4 & 5\\\hline y_i & 10 & 8 & 9 & 6 & 4 & 3\\\hline\end{array}$}\par}\smallskip Calcula el coeficiente de correlación. ¿Podemos afirmar que jugar a videojuegos hace bajar la nota?

\ex{2} Una heladería anota la temperatura máxima ($X$, en °C) y los helados vendidos ($Y$) durante seis días:\par\smallskip{\centering {\small$\begin{array}{|c|cccccc|}\hline x_i & 18 & 21 & 24 & 27 & 30 & 33\\\hline y_i & 35 & 52 & 58 & 75 & 84 & 102\\\hline\end{array}$}\par}\smallskip Calcula el coeficiente de correlación e interprétalo.

\ex{2} Calcula el coeficiente de correlación entre las notas de Matemáticas y de Física de la sección 4.

\ex{1} Si el coeficiente de correlación entre dos variables es $r=0{,}8$ y las desviaciones típicas son $\sigma_x=2$ y $\sigma_y=5$, ¿cuánto vale la covarianza?

\ex{1} Si $\sigma_{xy}=-12$, $\sigma_x=4$ y $\sigma_y=5$, ¿cuánto vale $r$? ¿Y si $\sigma_{xy}=3$, $\sigma_x=2$ y $\sigma_y=1$? ¿Es posible?

\ex{3} Una variable bidimensional tiene $N=10$, $\sum x_i=50$, $\sum y_i=80$, $\sum x_i^2=290$, $\sum y_i^2=730$ y $\sum x_iy_i=440$. Calcula $r$.


% =====================================================
\seccio{8}{Recta de regresión}

\exf{2} Con los datos del primer ejercicio de la sección 6, halla la recta de regresión de $Y$ sobre $X$ y la de $X$ sobre $Y$. Estima el valor de $y$ para $x=7$ y el de $x$ para $y=6$.\recorda{Recta de regresión de $Y$ sobre $X$: \ $y-\bar y=\dfrac{\sigma_{xy}}{\sigma_x^2}(x-\bar x)$. \ De $X$ sobre $Y$: \ $x-\bar x=\dfrac{\sigma_{xy}}{\sigma_y^2}(y-\bar y)$. Las estimaciones son fiables si $|r|$ está cerca de 1 y no nos salimos del intervalo de los datos.}

\ex{2} Con las alturas y los pesos de la sección 7, halla la recta de regresión del peso sobre la altura. ¿Qué peso se estima para una persona de 168 cm? ¿Y para un niño de 120 cm? ¿Son fiables las dos estimaciones?

\ex{2} Con los datos de la heladería de la sección 7, ¿cuántos helados se estima que se venderán un día con una máxima de 25 °C?

\ex{2} Con las sumas del último ejercicio de la sección 7 ($N=10$, $\sum x_i=50$, $\sum y_i=80$, $\sum x_i^2=290$, $\sum x_iy_i=440$), halla la recta de regresión de $Y$ sobre $X$ y estima $y$ para $x=6$.

\ex{3} Las rectas de regresión de una variable bidimensional son $y=0{,}5x+2$ (de $Y$ sobre $X$) y $x=1{,}2y-1$ (de $X$ sobre $Y$). Calcula $\bar x$, $\bar y$ y el coeficiente de correlación. (Pista: las dos rectas pasan por $(\bar x,\bar y)$, y el producto de sus pendientes es $\frac{\sigma_{xy}}{\sigma_x^2}\cdot\frac{\sigma_{xy}}{\sigma_y^2}=r^2$.)


% =====================================================
\seccio{9}{Otros tipos de relación. Correlación y causalidad}

\exf{2} Observa estas nubes de puntos. Describe la relación que hay entre las variables y di si el coeficiente de correlación lineal será próximo a 0 o no.\begin{tasks}(3)
  \task \par\begin{tikzpicture}[scale=0.22,baseline=(current bounding box.center)]\draw[muted!30,very thin] (0,0) grid (10.0,10.0);\draw[-latex,thick] (0,0)--(10.5,0) node[right,font=\scriptsize]{$X$};\draw[-latex,thick] (0,0)--(0,10.5) node[above,font=\scriptsize]{$Y$};\node[below,font=\tiny] at (0.0,0) {$0$};\node[below,font=\tiny] at (2.0,0) {$2$};\node[below,font=\tiny] at (4.0,0) {$4$};\node[below,font=\tiny] at (6.0,0) {$6$};\node[below,font=\tiny] at (8.0,0) {$8$};\node[below,font=\tiny] at (10.0,0) {$10$};\node[left,font=\tiny] at (0,0.0) {$0$};\node[left,font=\tiny] at (0,2.0) {$2$};\node[left,font=\tiny] at (0,4.0) {$4$};\node[left,font=\tiny] at (0,6.0) {$6$};\node[left,font=\tiny] at (0,8.0) {$8$};\node[left,font=\tiny] at (0,10.0) {$10$};\fill[accent] (0.3,1.416) circle (0.16);\fill[accent] (0.7,2.819) circle (0.16);\fill[accent] (1.1,4.031) circle (0.16);\fill[accent] (1.5,5.243) circle (0.16);\fill[accent] (1.9,5.924) circle (0.16);\fill[accent] (2.3,6.955) circle (0.16);\fill[accent] (2.7,6.583) circle (0.16);\fill[accent] (3.1,7.695) circle (0.16);\fill[accent] (3.5,8.745) circle (0.16);\fill[accent] (3.9,8.755) circle (0.16);\fill[accent] (4.3,9.31) circle (0.16);\fill[accent] (4.7,8.504) circle (0.16);\fill[accent] (5.1,8.959) circle (0.16);\fill[accent] (5.5,8.608) circle (0.16);\fill[accent] (5.9,8.769) circle (0.16);\fill[accent] (6.3,8.497) circle (0.16);\fill[accent] (6.7,7.404) circle (0.16);\fill[accent] (7.1,7.117) circle (0.16);\fill[accent] (7.5,6.548) circle (0.16);\fill[accent] (7.9,6.556) circle (0.16);\fill[accent] (8.3,5.507) circle (0.16);\fill[accent] (8.7,3.8) circle (0.16);\fill[accent] (9.1,3.473) circle (0.16);\fill[accent] (9.5,1.479) circle (0.16);\end{tikzpicture}
  \task \par\begin{tikzpicture}[scale=0.22,baseline=(current bounding box.center)]\draw[muted!30,very thin] (0,0) grid (10.0,10.0);\draw[-latex,thick] (0,0)--(10.5,0) node[right,font=\scriptsize]{$X$};\draw[-latex,thick] (0,0)--(0,10.5) node[above,font=\scriptsize]{$Y$};\node[below,font=\tiny] at (0.0,0) {$0$};\node[below,font=\tiny] at (2.0,0) {$2$};\node[below,font=\tiny] at (4.0,0) {$4$};\node[below,font=\tiny] at (6.0,0) {$6$};\node[below,font=\tiny] at (8.0,0) {$8$};\node[below,font=\tiny] at (10.0,0) {$10$};\node[left,font=\tiny] at (0,0.0) {$0$};\node[left,font=\tiny] at (0,2.0) {$2$};\node[left,font=\tiny] at (0,4.0) {$4$};\node[left,font=\tiny] at (0,6.0) {$6$};\node[left,font=\tiny] at (0,8.0) {$8$};\node[left,font=\tiny] at (0,10.0) {$10$};\fill[accent] (0.3,0.229) circle (0.16);\fill[accent] (0.7,0.015) circle (0.16);\fill[accent] (1.5,0.467) circle (0.16);\fill[accent] (1.9,0.188) circle (0.16);\fill[accent] (2.3,0.252) circle (0.16);\fill[accent] (2.7,0.707) circle (0.16);\fill[accent] (3.1,0.738) circle (0.16);\fill[accent] (3.5,0.547) circle (0.16);\fill[accent] (3.9,1.002) circle (0.16);\fill[accent] (4.3,0.933) circle (0.16);\fill[accent] (4.7,1.169) circle (0.16);\fill[accent] (5.1,1.13) circle (0.16);\fill[accent] (5.5,1.732) circle (0.16);\fill[accent] (5.9,1.883) circle (0.16);\fill[accent] (6.3,2.348) circle (0.16);\fill[accent] (6.7,2.698) circle (0.16);\fill[accent] (7.1,2.858) circle (0.16);\fill[accent] (7.5,3.431) circle (0.16);\fill[accent] (7.9,3.974) circle (0.16);\fill[accent] (8.3,4.721) circle (0.16);\fill[accent] (8.7,5.577) circle (0.16);\fill[accent] (9.1,6.755) circle (0.16);\fill[accent] (9.5,8.16) circle (0.16);\end{tikzpicture}
  \task \par\begin{tikzpicture}[scale=0.22,baseline=(current bounding box.center)]\draw[muted!30,very thin] (0,0) grid (10.0,10.0);\draw[-latex,thick] (0,0)--(10.5,0) node[right,font=\scriptsize]{$X$};\draw[-latex,thick] (0,0)--(0,10.5) node[above,font=\scriptsize]{$Y$};\node[below,font=\tiny] at (0.0,0) {$0$};\node[below,font=\tiny] at (2.0,0) {$2$};\node[below,font=\tiny] at (4.0,0) {$4$};\node[below,font=\tiny] at (6.0,0) {$6$};\node[below,font=\tiny] at (8.0,0) {$8$};\node[below,font=\tiny] at (10.0,0) {$10$};\node[left,font=\tiny] at (0,0.0) {$0$};\node[left,font=\tiny] at (0,2.0) {$2$};\node[left,font=\tiny] at (0,4.0) {$4$};\node[left,font=\tiny] at (0,6.0) {$6$};\node[left,font=\tiny] at (0,8.0) {$8$};\node[left,font=\tiny] at (0,10.0) {$10$};\fill[accent] (1.0,0.604) circle (0.16);\fill[accent] (1.15,1.414) circle (0.16);\fill[accent] (1.3,1.005) circle (0.16);\fill[accent] (1.45,0.972) circle (0.16);\fill[accent] (1.6,1.582) circle (0.16);\fill[accent] (1.75,1.177) circle (0.16);\fill[accent] (1.9,0.979) circle (0.16);\fill[accent] (2.05,1.177) circle (0.16);\fill[accent] (9.0,9.0) circle (0.16);\end{tikzpicture}
\end{tasks}

\ex{2} En un estudio de varias ciudades se ha encontrado una correlación positiva fuerte entre el número de bomberos que actúan en un incendio y los daños que causa. ¿Quiere decir que los bomberos causan los daños? ¿Qué explicación puede haber?

\ex{1} Razona si son verdaderas o falsas: \ a) Si $r=0$, las variables no tienen ninguna relación. \ b) Si $r=-0{,}98$, la relación es muy débil porque $r$ es negativo. \ c) Si todos los puntos están sobre una recta creciente, $r=1$.

\vfill
\begin{center}\color{muted}\small\faCheckCircle[regular]\ \ ¡Ánimo! Comprueba tus resultados con la calculadora bidimensional de los apuntes interactivos.\end{center}

\end{document}
